\tutorial{2}{Your own groundwater model}

\begin{goal}
You take an ordinary Raven model --- Tutorial Creek, a small catchment of two subbasins --- and give it a groundwater model
that Raven builds: aquifer classes and profiles, a DEM and a bedrock surface (GeoTIFF, and the same as ESRI ASCII grids),
input files in a map projection (\cmd{:InputCRS}), a grid and its refinement, head boundaries with and without a
\cmd{:BoundaryHead} series, a pumping well and observation wells. Then you run it and check the budget. Time: about an hour.
\end{goal}

\section{Tutorial Creek}
Set up the terminal as in Tutorial~1 (\cmd{PKG}, \cmd{RAVEN}, \cmd{TUT}, \cmd{TOOLS}) and copy the start state:
\begin{shell}
$ mkdir -p ~/work
$ cp -r $TUT/t2_own_model/start ~/work/t2
$ cd ~/work/t2
\end{shell}
\term{t2_ls_start.txt}
Tutorial Creek (Figure~\ref{fig:t2basin}) covers \V{t2:areaTotal}~km\textsuperscript{2} in UTM zone 11N. The creek runs west
to east along a valley floor 1.4~km wide; forested slopes rise about 150~m to the north and south; a bare-rock ridge
closes the upper subbasin to the north. Subbasin~1 (UPPER) ends in a small lake, CreekLake, which Raven routes as a
reservoir; subbasin~2 (LOWER) has a wetland and is gauged at its outlet. The \V{t2:nhru} HRUs:
\begin{center}\small
\begin{tabular}{rllrl}\toprule
HRU & subbasin & land & area [km\textsuperscript{2}] & groundwater (this tutorial)\\\midrule
1 & UPPER & valley floor & \V{t2:area1} & VALLEY\_FILL\\
2, 3 & UPPER & north and south slopes & \V{t2:area2}, \V{t2:area3} & SLOPE\_TILL\\
4 & UPPER & bare-rock ridge & \V{t2:area4} & none\\
5 & UPPER & CreekLake (reservoir) & \V{t2:area5} & VALLEY\_FILL\\
6 & LOWER & valley floor & \V{t2:area6} & VALLEY\_FILL\\
7, 8 & LOWER & north and south slopes & \V{t2:area7}, \V{t2:area8} & SLOPE\_TILL\\
9 & LOWER & wetland & \V{t2:area9} & VALLEY\_FILL\\\bottomrule
\end{tabular}\end{center}
\fig[0.86]{t2_basin.png}{Tutorial Creek: the HRUs (numbers), the creek, the lake (5) and the wetland (9), and the
groundwater features of this tutorial: a head boundary where the valley leaves the map, a pond with a specified level, a
pumping well and three observation wells.\label{fig:t2basin}}

The start folder is a complete Raven model \emph{without} groundwater, plus the data a groundwater model needs:
\begin{itemize}
\item \cmd{Creek.rvi}, \cmd{.rvh}, \cmd{.rvp}, \cmd{.rvc}, \cmd{.rvt} and \cmd{CreekMet.rvt}: the Raven model, with daily
      weather for 2010--2014. Its runs cover 2010--2012;
\item \cmd{hrus.geojson} (HRU polygons, field \cmd{HRU\_ID}) and \cmd{rivers.geojson} (the creek, field \cmd{SUBBASIN}), in UTM
      zone 11N (EPSG:32611);
\item \cmd{dem.tif} and \cmd{bedrock.tif}: land surface and bedrock, 50~m GeoTIFF rasters in the same projection;
      \cmd{dem.asc} and \cmd{bedrock.asc} hold the same values as ESRI ASCII grids;
\item \cmd{outlet.geojson} (a line across the valley at its eastern edge) and \cmd{pond.geojson} (a gravel-pit pond);
\item \cmd{CreekGW.rvt}: the pumping record of well PW\_1 and the measured level of the pond.
\end{itemize}

\subsection{Run the model without groundwater first}
\term{t2_run_nogw.txt}
The model runs in a fraction of a second; the mean flow at the outlet is \V{t2:qmeanNoGW}~m\textsuperscript{3}/s. Its
deep store is a linear reservoir: \cmd{Creek.rvi} sends water from the subsoil to \cmd{DEEP\_STORE} (\cmd{SOIL[2]}), which
drains to the creek.
\needspace{30\baselineskip}
\showfile{t2_own_model/start/Creek.rvi}
You will replace \cmd{DEEP\_STORE} by a MODFLOW~6 model under every HRU except the ridge, in four edits of the Raven
files and one new file, \cmd{Creek.rvg}.

\section{Which HRUs have groundwater (\texttt{.rvh})}
The \cmd{AQUIFER\_PROFILE} column of the HRU table names the aquifer profile under each HRU (\cmd{[NONE]}: no groundwater).
The valley floor, the lake and the wetland lie on the valley fill, the slopes on till; the ridge keeps \cmd{[NONE]}. An HRU
group \cmd{AquiferHRUs} lists the coupled HRUs, for the process conditions of the \cmd{.rvi}. Edit \cmd{Creek.rvh}; the exact
change is (lines starting with \cmd{-} are removed, \cmd{+} added):
\diff[basicstyle=\ttfamily\tiny]{t2_rvh.diff}

\section{The hydrogeology (\texttt{.rvp})}
Add three tables to the end of \cmd{Creek.rvp}:
\diff{t2_rvp.diff}
\begin{itemize}
\item \cmd{:AquiferClasses}: the materials. \cmd{K\_HORIZ} and \cmd{K\_VERT} are hydraulic conductivities [m/d],
      \cmd{SPEC\_STORAGE} the specific storage [1/m], \cmd{SPEC\_YIELD} the specific yield; \cmd{POROSITY} is used by the
      unsaturated columns of Tutorial~3 and by particle tracking (Tutorial~5).
\item \cmd{:AquiferProfiles}: the layer stack under an HRU, from the top: class, thickness (or \cmd{TO\_BEDROCK}: down to the
      bedrock raster) and layer type. Both profiles have two layers, so layer~1 of the valley joins layer~1 of the slopes
      where they meet. \cmd{AQUIFER} layers are convertible (they can dry and rewet); \cmd{CONFINED\_AQUIFER},
      \cmd{AQUITARD} and \cmd{AQUICLUDE} are the other types (user guide, section~4.2).
\item \cmd{:AquiferProfileParameters}: the initial water-table depth and the conductance of the seepage face where the water
      table reaches the land surface. Other parameters (river-bed conductivity, water-table ET, recharge delay, \dots)
      keep their defaults.
\end{itemize}

\section{Switch the coupling on and send recharge to it (\texttt{.rvi})}
\diff{t2_rvi.diff}
\cmd{:GroundwaterModel MODFLOW6} switches the coupling on. In the coupled HRUs the subsoil now percolates to
\cmd{GROUNDWATER}, which Raven passes to MODFLOW as recharge every day; the ridge keeps its \cmd{DEEP\_STORE}. Groundwater
that seeps out at the surface returns to the lowest soil layer (\cmd{DEEP\_STORE}), whose baseflow process takes it to the
creek; exchange with the river cells goes straight to the reaches.
\begin{trouble}
Raven stops if a process takes water \emph{out} of \cmd{GROUNDWATER} in a coupled HRU (MODFLOW owns that storage), if no
process sends water to \cmd{GROUNDWATER}, or if a process sends water to \cmd{GROUNDWATER} in an HRU without an
\cmd{AQUIFER\_PROFILE} (``\dots which has no AQUIFER\_PROFILE''): restrict such processes with \cmd{:-->Conditional}, as here.
\end{trouble}

\section{The groundwater file (\texttt{.rvg})}
Create \cmd{Creek.rvg} with a text editor; its complete content:
\showfile{t2_own_model/finished/Creek.rvg}
Line by line:

\subsection{Input files and their projection}
\cmd{:InputCRS EPSG:32611} says that \emph{every} GeoJSON file and raster of the \cmd{.rvg} is in UTM zone 11N. Without it,
Raven expects longitude and latitude (EPSG:4326). Raven checks this: a GeoJSON file whose \cmd{"crs"} member names another
projection, a GeoTIFF whose EPSG code differs, or coordinates that look like longitude/latitude stop the run with the
reason. Outputs (\cmd{grid\_cells.geojson}, the \cmd{.rvh} coordinates) stay in longitude/latitude.

\cmd{:HRUGeometry} reads the HRU polygons and the name of the field that holds the HRU ID, \cmd{:RiverGeometry} the creek
lines and the field that holds their subbasin (without a field, each piece goes to the subbasin of the HRU that dominates
its cell). \cmd{:LandSurface} and \cmd{:BedrockSurface} read the rasters; each cell takes the mean of the raster over its
area. Raven reads GeoTIFF files itself (one band, uncompressed, LZW, DEFLATE or PackBits, north-up) and ESRI ASCII grids;
no GIS software is needed.

\subsection{Grid and start}
\cmd{:GridCellSize 250} lays a regular grid of 250-m cells over the coupled HRUs; a column is active where coupled HRUs
cover at least a quarter of it (\cmd{:MinCellCoverage}, default 0.25). \cmd{:GWInitialization STEADY\_STATE 1.2} starts the
run from the steady state of a uniform recharge of 1.2~mm/d instead of the profiles' \cmd{INITIAL\_HEAD\_DEPTH}, so the
first months are not spent filling or draining the aquifer.

\subsection{Wells and observation wells}
With \cmd{:InputCRS}, wells are placed by \cmd{X}, \cmd{Y} in the projection (otherwise \cmd{LATITUDE}, \cmd{LONGITUDE}).
PW\_1 is screened between 958 and 945~m, in the sand of layer~2. Its rate comes from a \cmd{:WellFlowRate} series in the
\cmd{.rvt} file (negative: pumping); MODFLOW reduces it automatically if the cell runs dry (\cmd{wells actual} in
\cmd{GWBudget.csv}). Each observation well gives the elevation of its screen; Raven writes the head there every day to
\cmd{GWHeads.csv}, and observed heads (\cmd{:ObservationData GW\_HEAD}, Tutorial~6) are scored in \cmd{Diagnostics.csv}.

\subsection{Head boundaries, with and without a series}
\begin{itemize}
\item \cmd{:GeneralHeadBoundary OUTLET outlet.geojson HEAD 958.0 CONDUCTANCE 500.0 LAYERS ALL}: groundwater leaves the valley
      through its eastern edge towards a head of 958~m. The line selects the cells it crosses; the conductance is per cell
      [m\textsuperscript{2}/d]. Without a \cmd{:BoundaryHead} series the head stays at 958~m.
\item \cmd{:SpecifiedHeadBoundary POND pond.geojson HEAD 976.5 LAYERS TOP}: the cells whose centre lies in the pond polygon
      are held at the pond level. \cmd{CreekGW.rvt} has a \cmd{:BoundaryHead POND} series (the measured level, \V{t2:pondmin} to
      \V{t2:pondmax}~m over the year); it replaces \cmd{HEAD} day by day.
\end{itemize}
A boundary is placed only in layers whose bottom lies below its head (for a series: below its highest value); on a day
when the head lies below a placed cell's bottom the cell is switched off (counted in \cmd{GWModelSummary.txt}).

\section{The well and pond series (\texttt{.rvt})}
\cmd{CreekGW.rvt} is ready in the folder:
\term{t2_gwrvt_head.txt}
\cmd{:WellFlowRate 1 m3/d} refers to the well with ID~1 of \cmd{:Wells}; \cmd{:BoundaryHead POND m} to the boundary named
POND. Both are daily series (start, interval 1~day, number of values). Tell Raven to read the file:
\diff{t2_rvt.diff}

\section{Run}
\term{t2_run.txt}
The run takes about \V{t2:runsec}~seconds. \cmd{Groundwater steady-state initialization converged in \V{t2:ssiter} outer
iterations} confirms the steady start.
\begin{tocheck}
\cmd{Exiting Gracefully: Successful Simulation}, and \cmd{out/Raven\_errors.txt} without \cmd{ERROR} lines.
\end{tocheck}
\begin{trouble}
\begin{itemize}
\item \emph{``HRU \dots has an AQUIFER\_PROFILE but no polygon''}: the IDs in \cmd{hrus.geojson} and \cmd{Creek.rvh} do not
      match, or the ID field is not \cmd{HRU\_ID}.
\item \emph{``\dots does not overlap any active groundwater cell''}: a small HRU under large cells; reduce
      \cmd{:GridCellSize} or \cmd{:MinCellCoverage}.
\item \emph{``:WellFlowRate 1 in the .rvt file names no well''} or \emph{``:BoundaryHead POND \dots names no
      :GeneralHeadBoundary or :SpecifiedHeadBoundary''}: the ID or name differs between the \cmd{.rvt} and the \cmd{.rvg}.
\item A forgotten \cmd{:InputCRS} --- see Section~\ref{sec:t2crs}.
\end{itemize}
\end{trouble}

\subsection{What Raven built}
\term{t2_summary.txt}
\begin{tocheck}
\begin{itemize}
\item The grid: \V{t2:nlay}~layers of \V{t2:nrow}$\times$\V{t2:ncol} cells; \V{t2:activecols} active columns (the
      \V{t2:inactivecols} others lie under the ridge or outside the HRUs); \V{t2:fidelity}\% of the active area lies in
      each cell's dominant HRU.
\item \V{t2:rivcells} river cells, one well cell, \V{t2:ghbcells} GHB cells (in both layers) and \V{t2:chdcells} CHD
      cells under the pond.
\item The observation wells: OW\_NEAR in layer \V{t2:owlay101}, OW\_UPPER in layer \V{t2:owlay102}, OW\_SLOPE in layer
      \V{t2:owlay103} (the top layer is dry there).
\item The HRU table: \cmd{linked\_fraction\_of\_polygon} close to 1 for every HRU. A low value means an HRU outside the grid
      or a wrong projection.
\end{itemize}
\end{tocheck}
Figure~\ref{fig:t2grid} shows the active cells (from \cmd{out/mf6/grid\_cells.geojson}) over the HRUs.
\fig[0.86]{t2_grid.png}{The active cells of the grid Raven built, over the HRUs: none under the ridge (HRU~4), which has no
aquifer.\label{fig:t2grid}}

\subsection{Check the budget}
\term{t2_budget_head.txt}
Over the three years the mean budget is (m\textsuperscript{3}/d):
\begin{center}\small
\begin{tabular}{lr@{\qquad}lr}\toprule
into the aquifer & & out of the aquifer &\\\midrule
recharge from Raven & \V{t2:rech} & seepage to Raven & \V{t2:seep}\\
river leakage & \V{t2:fromriv} & discharge to the creek & \V{t2:toriv}\\
storage release & \V{t2:stor} & GHB OUTLET & \V{t2:ghb}\\
 & & well PW\_1 & \V{t2:wel}\\
 & & pond (CHD) & \V{t2:chd}\\\bottomrule
\end{tabular}\end{center}
The recharge is \V{t2:rechmm}~mm/year over the coupled HRUs; \V{t2:seeppct}\% of it seeps out again, mostly at the foot of
the slopes and on the valley floor. MODFLOW converged on every day (\V{t2:unconv} unconverged steps, at most
\V{t2:maxouter} outer iterations) and its largest daily balance error is \V{t2:balerr}~\% of the flow. The audit tool of the
Liard example works for any coupled run:
\term{t2_audit.txt}
\begin{tocheck}
The recharge identity is 0, the cumulative MODFLOW residual \V{t2:auditRes} of the gross flow, no unconverged steps, the
well delivered what was asked (\cmd{wells delivered / specified} 1), Raven's own balance error is \V{t2:auditMB}~mm
(\cmd{Raven MB error max}) and \cmd{ERROR lines} 0.
\end{tocheck}

\subsection{Look at the results}
\term{t2_plot.txt}
\fig{t2_gw_budget.png}{The budget of Tutorial Creek (7-day means): winter recharge, seepage and discharge to the creek
through the year, the outflow at the GHB, the well and the pond.}
\fig{t2_gw_watertable.png}{Water-table depth: about \V{t2:wt6}~m under the valley floor (HRU~6) and \V{t2:wt7}~m under the
north slope of the lower subbasin (HRU~7).}
\fig[0.75]{t2_gw_map.png}{Water-table depth at the last saved head field: shallow on the valley floor, tens of metres under
the upper slopes.}
\fig{t2_gw_obswells.png}{Heads at the three observation wells (\cmd{GWHeads.csv}). OW\_NEAR, 200~m from the pumping well,
follows its seasonal pumping.}
Groundwater changes the streamflow too (Figure~\ref{fig:t2hyd}): the mean flow drops from \V{t2:qmeanNoGW} to
\V{t2:qmean}~m\textsuperscript{3}/s, about \V{t2:qdrop}~m\textsuperscript{3}/d. Most of that water leaves the basin
underground, \V{t2:ghb}~m\textsuperscript{3}/d at the GHB and \V{t2:wel}~m\textsuperscript{3}/d through the well. The rest
is held in Raven: the seepage from the aquifer collects in \cmd{DEEP\_STORE}, which gains \V{t2:deepgw}~mm over the three
years with groundwater against \V{t2:deepnogw}~mm without (about \V{t2:deeprate}~m\textsuperscript{3}/d over the basin) ---
the store spins up to its new inflow. The aquifer itself releases water from storage (\V{t2:stor}~m\textsuperscript{3}/d in
the table above). The model is not calibrated: Tutorial~6 does that.
\fig{t2_hydrograph.png}{Streamflow at the outlet in 2011, without and with the groundwater model.\label{fig:t2hyd}}

\section{Variations}
Try each change, run, and undo it before the next one.

\subsection{ESRI ASCII grids instead of GeoTIFF}
\diff{t2_ascii.diff}
\begin{shell}
$ $RAVEN Creek -o out_ascii/
\end{shell}
The two pairs of rasters hold the same values (the GeoTIFF files as 32-bit floats), so the heads in
\cmd{out\_ascii/GWHRUState.csv} differ from those of \cmd{out/} by at most \V{t2:asciidiff}~m. An ESRI ASCII grid has no
projection of its own: \cmd{:InputCRS} applies to it, and Raven checks only that its extent lies over the HRUs.

\subsection{Forgetting \texttt{:InputCRS}}\label{sec:t2crs}
\diff{t2_nocrs.diff}
\term{t2_nocrs_err.txt}
Raven stops before building anything: without \cmd{:InputCRS} the coordinates of the well (and of every input file) would
be read as longitude and latitude. The message says what to do. (Had the wells been given by latitude and longitude, the
next message would have named \cmd{hrus.geojson}, whose \cmd{"crs"} member says EPSG:32611.)

\subsection{Finer cells along the creek}
\diff{t2_refine.diff}
\begin{shell}
$ $RAVEN Creek -o out_refined/
\end{shell}
\term{t2_refine_summary.txt}
A regular grid is refined in whole rows and columns: \V{t2:refnrow}$\times$\V{t2:refncol} cells and \V{t2:refactive}
active columns instead of \V{t2:nrow}$\times$\V{t2:ncol} and \V{t2:activecols}, \V{t2:refriv} river cells instead of
\V{t2:rivcells}, and about \V{t2:refsec}~seconds per run instead of \V{t2:runsec}. The finer cells along the creek change
the discharge to it (\V{t2:reftoriv} instead of \V{t2:toriv}~m\textsuperscript{3}/d): the narrower river cells take
less of the groundwater, and more of it seeps out into the soil next to them instead (seepage to Raven \V{t2:refseep}
instead of \V{t2:seep}~m\textsuperscript{3}/d). The mean water-table depth of HRU~6 becomes \V{t2:refwt6}~m. \cmd{:GridType QUADTREE} refines locally instead (user guide, section~4.3a).

\section{The finished state}
\cmd{\$TUT/t2\_own\_model/finished} holds every file as it is at the end of this tutorial (without the variations). Compare
with yours: \cmd{diff -r \$TUT/t2\_own\_model/finished \textasciitilde/work/t2} lists only your output folders.

\section{Adapting this to your own model}
\begin{enumerate}
\item Export HRU polygons with the HRU IDs of your \cmd{.rvh} (GeoJSON), the river lines with their subbasin, a DEM and an
      aquifer-base raster; all in longitude/latitude or all in one projection (\cmd{:InputCRS}).
\item Choose aquifer profiles; fill \cmd{AQUIFER\_PROFILE}; define the classes and profiles in the \cmd{.rvp}.
\item Send the water your model sends to deep groundwater into \cmd{GROUNDWATER} for the coupled HRUs, and make sure nothing
      draws from it there.
\item Write the \cmd{.rvg}; start with cells of a few hundred metres to a few kilometres, then refine.
\item Run; read \cmd{GWModelSummary.txt}, run the audit, plot. Then calibrate (Tutorial~6).
\end{enumerate}
